\documentclass{article}
\usepackage[T1]{fontenc}
\usepackage{lmodern}
\usepackage{graphicx}
\usepackage{amsmath,amssymb}
\usepackage[a4paper,margin=2cm]{geometry}
\usepackage[absolute,overlay]{textpos}
\setlength{\TPHorizModule}{1mm}
\setlength{\TPVertModule}{1mm}
\usepackage[indonesian]{babel}
\usepackage{hyperref}
\setlength{\emergencystretch}{2em}
% unit: Halaman 101

\begin{document}

% === Halaman 101 (llm) ===

Periksa bahwa:

\begin{align*}
\mathbf{KK}^{-1} &= \begin{pmatrix} 3 & 10 \\ 15 & 9 \end{pmatrix} \begin{pmatrix} 5 & 6 \\ 9 & 19 \end{pmatrix} = \begin{pmatrix} 3\cdot 5 + 10\cdot 9 & 3\cdot 6 + 10\cdot 19 \\ 15\cdot 5 + 9\cdot 9 & 15\cdot 6 + 9\cdot 19 \end{pmatrix} \\ &= \begin{pmatrix} 105 & 208 \\ 156 & 261 \end{pmatrix} \pmod{26} = \begin{pmatrix} 1 & 0 \\ 0 & 1 \end{pmatrix}
\end{align*}

\end{document}
